% !TeX root = ../../main.tex

\lstset{
    basicstyle=\ttfamily\footnotesize,
    breaklines=true,
    frame=single,
    showstringspaces=false,
    columns=fullflexible,
    literate=
        {á}{{\'a}}1
        {é}{{\'e}}1
        {í}{{\'i}}1
        {ó}{{\'o}}1
        {ú}{{\'u}}1
        {ü}{{\"u}}1
        {ñ}{{\~n}}1
}

\section{Introducción}

En esta actividad se implementó un contador de palabras mediante
MapReduce. Para realizar la práctica se utilizó Hadoop 3.5.0 dentro de
contenedores Docker. El programa fue escrito en Python y se ejecutó
mediante Hadoop Streaming.

\section{Configuración de Hadoop}

El entorno se creó con Docker Compose y está formado por los siguientes
servicios:

\begin{itemize}
    \item NameNode: administra el sistema de archivos HDFS.
    \item DataNode: almacena los bloques de datos.
    \item ResourceManager: administra los recursos del clúster.
    \item NodeManager: ejecuta las tareas de MapReduce.
\end{itemize}

Para iniciar los servicios se utilizó el siguiente comando:

\begin{lstlisting}[language={}, caption={Inicio de los servicios}]
docker compose up -d
\end{lstlisting}

La ejecución de \texttt{docker compose ps} permitió comprobar que los
cuatro contenedores estaban activos.

\begin{figure}[H]
    \centering
    \includegraphics[width=0.95\textwidth]
    {\actividadactual/img/captura-4.png}
    \caption{Contenedores de Hadoop en ejecución.}
\end{figure}

\section{Implementación del contador de palabras}

El archivo de entrada contiene varias oraciones relacionadas con
Hadoop, Docker y MapReduce.

\lstinputlisting[
    language={},
    caption={Archivo de entrada}
]{\actividadactual/docker-hadoop/src/entrada.txt}

\subsection{Mapper}

El programa \texttt{mapper.py} recibe el texto por la entrada estándar,
convierte las palabras a minúsculas y genera un par formado por la
palabra y el número uno.

\lstinputlisting[
    language=Python,
    caption={Código del mapper}
]{\actividadactual/docker-hadoop/src/mapper.py}

Por ejemplo, si aparece dos veces la palabra \texttt{hadoop}, el mapper
genera inicialmente:

\begin{lstlisting}[language={}]
hadoop    1
hadoop    1
\end{lstlisting}

\subsection{Reducer}

El programa \texttt{reducer.py} recibe las palabras ordenadas y suma los
valores correspondientes a cada una. Su salida contiene cada palabra
junto con su número total de apariciones.

\lstinputlisting[
    language=Python,
    caption={Código del reducer}
]{\actividadactual/docker-hadoop/src/reducer.py}

\section{Ejecución en Hadoop}

Primero se creó el directorio de entrada en HDFS y se copió el archivo
de texto:

\begin{lstlisting}[language={}, caption={Carga del archivo en HDFS}]
docker compose exec -T namenode \
hdfs dfs -mkdir -p /user/hadoop/input

docker compose exec -T namenode \
hdfs dfs -put -f /workspace/entrada.txt \
/user/hadoop/input/
\end{lstlisting}

Después se ejecutó el trabajo mediante Hadoop Streaming:

\begin{lstlisting}[language={}, caption={Ejecución de Hadoop Streaming}]
docker compose exec -T namenode \
hadoop jar \
/opt/hadoop/share/hadoop/tools/lib/hadoop-streaming-3.5.0.jar \
-D mapreduce.job.name=conteo-palabras-python \
-files /workspace/mapper.py,/workspace/reducer.py \
-mapper "python3 mapper.py" \
-reducer "python3 reducer.py" \
-input /user/hadoop/input \
-output /user/hadoop/output
\end{lstlisting}

La ejecución terminó correctamente con el mapper y el reducer al
100\,\%.

\begin{figure}[H]
    \centering
    \includegraphics[width=0.95\textwidth]
    {\actividadactual/img/captura-1.png}
    \caption{Ejecución correcta del trabajo MapReduce.}
\end{figure}

También se verificó el trabajo desde la interfaz de YARN. La aplicación
apareció con estado \texttt{FINISHED} y resultado
\texttt{SUCCEEDED}.

\begin{figure}[H]
    \centering
    \includegraphics[width=0.95\textwidth]
    {\actividadactual/img/captura-3.png}
    \caption{Trabajo completado correctamente en YARN.}
\end{figure}

\section{Resultados}

Para consultar el resultado almacenado en HDFS se utilizó el siguiente
comando:

\begin{lstlisting}[language={}, caption={Consulta del resultado}]
docker compose exec -T namenode \
hdfs dfs -cat /user/hadoop/output/part-00000
\end{lstlisting}

El programa contó correctamente las palabras del archivo. Algunos de
los resultados obtenidos fueron:

\begin{center}
\begin{tabular}{lr}
    \toprule
    \textbf{Palabra} & \textbf{Cantidad} \\
    \midrule
    hadoop    & 4 \\
    y         & 3 \\
    datos     & 2 \\
    map       & 2 \\
    mapreduce & 2 \\
    procesa   & 2 \\
    reduce    & 2 \\
    \bottomrule
\end{tabular}
\end{center}

\begin{figure}[H]
    \centering
    \includegraphics[width=0.70\textwidth]
    {\actividadactual/img/captura-2.png}
    \caption{Resultado del conteo de palabras almacenado en HDFS.}
\end{figure}

\section{Conclusión}

La práctica permitió ejecutar un proceso MapReduce en un entorno Hadoop
instalado con Docker. El mapper generó los pares de palabra y valor,
mientras que el reducer sumó las apariciones. Finalmente, Hadoop
almacenó el resultado en HDFS y YARN confirmó que el trabajo terminó
correctamente.